\section{Motivating Example}

Consider the expression
\[
  (a \mid aa)^{*} b
\]
on inputs of the form \(a^n b\) and \(a^n c\).  A backtracking engine explores
ways of decomposing the prefix \(a^n\) into repetitions of \(a\) and \(aa\).
On \(a^n b\), some decompositions lead to a successful match; on \(a^n c\),
the engine may explore many decompositions before discovering that the final
literal cannot be matched.  This is the familiar shape of a ReDoS example:
the engine's trace tree grows rapidly with \(n\), while the input itself grows
linearly.

The automata view says that each decomposition is an accepting path in a
nondeterministic automaton for the expression prefix.  The engine view says that
each decomposition is a branch in an ordered backtracking tree.  The same
combinatorial object is therefore seen twice: once as an unordered set of NFA
runs, and once as a priority-ordered trace.  A proof relating the two views lets
us use automata-theoretic ambiguity to reason about the successful frontier of
the trace tree.

This example also illustrates a limitation.  ReDoS is not only about successful
matches.  Failing branches can dominate runtime, especially on near-miss inputs.
Our bridge is therefore a first layer: it relates successful trace leaves to
accepting runs.  A complete ReDoS account must also quantify failing prefixes,
which we discuss in Section~\ref{sec:redos}.
